\chapteropen{03}

\begin{cele}
\item policzyć, ile minut dziennie naprawdę masz, zanim cokolwiek zmienisz;
\item odróżnić mikroczas od stałego okna i wiedzieć, do czego nadaje się każdy z nich;
\item zamienić ekran na książkę tak, żeby nie polegało to na silnej woli;
\item przeliczyć koszt rezygnacji z jednej konkretnej rzeczy i podjąć tę decyzję świadomie.
\end{cele}

Większość osób, które mówią \glqq nie mam czasu\grqq{}, nigdy nie policzyła, ile go ma.
Zdanie brzmi jak diagnoza, a jest raczej hasłem. Zanim zaczniesz szukać godzin, potrzebujesz
liczby, a liczba zwykle okazuje się inna, niż podpowiada pamięć. Ten rozdział jest więc
rachunkiem: ile minut mieści się w dniu, gdzie one są, co kosztuje ich odzyskanie i co
z tego wynika dla roku. Liczby z przykładów są wyliczone w tekście, a nie zmierzone u kogoś.
Tam, gdzie chodzi o świat poza książką, a nie o własną arytmetykę, zaznaczam to osobno.

\section{Najpierw audyt, potem plan}

Przez siedem dni zapisuj każdy blok dłuższy niż pięć minut: co robiłeś, ile to trwało i czy
w tym czasie był w ręku ekran. Nie oceniaj. Interesuje cię inwentarz, a nie wyrzuty sumienia.
Po tygodniu zsumuj każdą kategorię osobno, bo pojedyncze bloki wyglądają niewinnie, a suma już nie.

Poniższa tabela jest przykładem, a nie danymi. Pokazuje, jak wygląda czytelny zapis dnia:

\begin{center}
\begin{tabular}{lrc}
\toprule
Blok & Minuty dziennie & Ekran w ręku \\
\midrule
Dojazd do pracy i z powrotem (2 $\times$ 25 min) & 50 & tak \\
Kolejki, oczekiwanie w urzędzie, na wizytę & 18 & tak \\
Przerwa w pracy & 20 & tak \\
Wieczór przed snem & 45 & tak \\
Obiad & 25 & nie \\
\bottomrule
\end{tabular}
\end{center}

Zwróć uwagę na ostatnią kolumnę. Ona mówi więcej niż same minuty, bo wskazuje, które bloki
już są czytelne, a które zajmuje urządzenie, które trzeba będzie odłożyć. Blok z ekranem
nie jest stracony. Jest zajęty czymś, co można zamienić, i to właśnie ta zamiana
będzie tematem sekcji czwartej.

Audyt ma jeszcze jedną funkcję: odsiewa mity. Wiele osób jest przekonanych, że cały ich dzień
jest wypełniony. Tydzień zapisów zwykle pokazuje, że wypełniony jest mniej więcej co drugi
kwadrans, a reszta rozchodzi się na czekanie, przełączanie i przeglądanie.

\section{Mikroczas: kolejki, dojazd, przerwy}

Mikroczas to bloki od pięciu do trzydziestu minut, których nie planujesz w kalendarzu, a które
pojawiają się same. Mają jedną zaletę: powtarzają się codziennie. Mają też wadę: są krótkie
i łatwo je zmarnować.

Policzmy przykładowy dzień z tabeli. Dojazd to 50 minut, kolejki i oczekiwanie to 18, przerwa
w pracy to 20. Razem daje to \textbf{88 minut} dziennie. Nie wszystkie te minuty nadają się do
czytania. W tłoku, w kolejce na zewnątrz albo w przerwie, którą spędzasz z kolegami przy stole,
książka nie wchodzi w grę. Załóżmy ostrożnie, że wykorzystasz jedną trzecią: $88 / 3 \approx 29$
minut dziennie.

Teraz zestawmy to z liczbą z tablicy minut. Dla dwunastu książek rocznie potrzeba
\val{need.12} minuty dziennie, a dla dwudziestu czterech \val{need.24} minut. Z mikroczasu
w tym przykładzie wychodzi około 29 minut, czyli wystarczy na dwadzieścia cztery książki
i zostaje jeszcze zapas. Nie trzeba do tego żadnego nowego okna w kalendarzu. Wystarczy,
że przestaniesz traktować te minuty jako nieistniejące.

Dojazd wymaga jednak osobnej uwagi. Jeśli jedziesz autobusem lub pociągiem i masz miejsce,
możesz czytać. Jeśli prowadzisz samochód, nie czytasz, a słuchanie to inna czynność, o której
jest rozdział 5. Każdy blok ma więc własny zakres: to, co jest realne w danej sytuacji,
a nie to, co wyglądałoby dobrze na wykresie.

\begin{falle}
\textbf{\glqq Nie mam czasu, bo nie mam półtorej godziny\grqq{}.} Ten argument zakłada, że
czytanie wymaga jednego długiego bloku. Dla lekkiej książki faktycznej i dla powieści tak bywa
rzadziej, niż się wydaje. Dla trudnego tekstu jest odwrotnie: przerywanie wątku kosztuje, a każde
wznowienie zaczyna się od przypominania sobie, na czym się skończyło. Mikroczas dobrze działa
na tekst, który da się podjąć w dziesięciu minutach. Dla wymagającej lektury potrzebujesz
stałego okna, o którym jest następna sekcja. Mikroczas nie jest rozwiązaniem uniwersalnym.
\end{falle}

\section{Stałe okna: rano 20 minut, wieczór 30}

Mikroczas jest nieprzewidywalny, więc dobrze jest mieć też stałe okna, czyli bloki, które zawsze
są zarezerwowane. Dwa z nich wystarczą, bo każde ma inną funkcję. Rano 20 minut to okno na
czytanie przed tym, co może zjeść całą resztę dnia. Wieczór 30 minut to okno na lekturę
wymagającą większej koncentracji.

Oba okna zapisz w kalendarzu jak spotkania. Brzmi to banalnie, ale tylko zapisane okno przetrwa
pierwszy zły dzień. Niezapisane znika w pierwszym tygodniu, bo ktoś zaprosi cię na kolację albo
przyjdzie awaria.

Rachunek rocznej wydajności wygląda tak. Zakładamy, że czytasz z tempem 250 słów na minutę
i że przeciętna książka ma \val{typical.words} słów. \verify{typowe tempo czytania ciszą wynosi
około 250 słów na minutę} --- to założenie, które warto sprawdzić na własnym tekście, mierząc
czas na jednej stronie. Przy takim założeniu książka to około \val{typical.hours} godziny czytania.
Okno poranne, 20 minut dziennie, daje \val{books.20.250} książek rocznie. Okno wieczorne,
30 minut dziennie, daje \val{books.30.250} książek. Razem, przy tych samych założeniach,
wychodzi 57 książek w roku.

Ta liczba jest górną granicą, a nie prognozą. Dni wypadają, teksty bywają trudniejsze niż
przeciętne, a tempo spada, gdy czytasz coś, co wymaga zatrzymywania się przy każdym akapicie.
Traktuj ją jak sufit, który mówi, co jest możliwe, a nie co się stanie. Realne liczby wyjdą z
własnego tygodnia, a nie z tabeli.

\section{Ekran, ładowarka i scrollowanie}

Zamiana ekranu na książkę jest najtrudniejszą częścią tego rozdziału, bo ekran nie czeka na
decyzję. Działa tam, gdzie jest, a jego rozpraszanie wymaga zero wysiłku. Dlatego najskuteczniejsze
zmiany dotyczą miejsca, a nie charakteru.

Najprostsza z nich: ładowarkę przenieś poza sypialnię. Telefon ładuje się w kuchni albo w przedpokoju,
a przy łóżku leży książka. Budzik kup osobny, jeśli dotąd budziłeś się telefonem. Nie chodzi o
to, że telefon jest zły. Chodzi o to, że przy łóżku wieczorne przeglądanie zaczyna się samo,
a sen może na tym tracić. Dorośli potrzebują \verify{dorośli zwykle potrzebują 7--9 godzin snu}
i wszystko, co zabiera z tej puli, trzeba liczyć jako koszt, a nie jako drobnostkę.

Teraz rachunek. Załóżmy, że scrollowanie zajmuje ci 45 minut dziennie, w sumie za cały rok:
\[ 45 \times 365 = 16\,425 \text{ minut} = 273{,}75 \text{ godziny}. \]
To jest ponad jedenaście dób. Przy założeniach z poprzedniej sekcji oznacza to
\val{books.45.250} książek, czyli tyle samo, ile dałoby codzienne czytanie przez 45 minut.
Nie musisz rezygnować z przeglądania, żeby to zobaczyć. Wystarczy, że wiesz, ile kosztuje.

Zamiana nie jest jednak wymianą jeden do jednego. Część scrollowania to odpoczynek, którego
nie chcesz zastąpić wysiłkiem, i to jest w porządku. Zamień tę część, która jest wyłącznie
odruchem: otwieranie aplikacji w przerwie na jedną rzecz, która potrafi zabrać pół godziny.

\section{Rezygnacja z czegoś konkretnego}

Czas na czytanie trzeba czasem wziąć od czegoś, a nie tylko znaleźć w luce. Najłatwiej wybrać
jedną konkretną rzecz i policzyć ją w tygodniu. Jeśli sprawdzanie wiadomości w jednej aplikacji
zajmuje ci 20 minut dziennie, to:
\[ 20 \times 7 = 140 \text{ minut tygodniowo} = 2 \text{ godziny } 20 \text{ minut}, \]
a w roku:
\[ 20 \times 365 = 7\,300 \text{ minut} \approx 121{,}7 \text{ godziny}. \]
Przy tempie z sekcji trzeciej to 140 minut tygodniowo, czyli około 35\,000 słów, a więc niemal
połowa przeciętnej książki w tygodniu. Nie musisz z tego rezygnować całkowicie. Wystarczy
decyzja o jednej zmianie: wyłączone powiadomienia, usunięta aplikacja z ekranu głównego, ograniczenie
do jednego razu dziennie. Rezygnacja z czegoś konkretnego działa lepiej niż postanowienie, że
będziesz mniej przeglądać, bo postanowienie nie ma miejsca w kalendarzu.

Podobnie jest z podróżami i urlopem. Podróż pociągiem trwająca sześć godzin to 360 minut. Jeśli
przeczytasz z niej 60\%, zyskujesz 216 minut, czyli 3 godziny 36 minut lektury z jednego dnia.
To dużo, a jednak wiele osób traktuje takie godziny jako stracone, bo nie są w planie. Urlop
ma osobny rachunek. Dwa tygodnie po dwie godziny dziennie to 28 godzin, czyli więcej niż
czytanie przez cały tydzień przy oknie porannym. Urlop nie zwiększa jednak tempa sam z siebie.
Jeśli czytasz tylko tam, nawyk zwykle nie wraca wraz z powrotem do pracy. Zabierz na urlop tę
samą książkę, która ma być czytana w domu, zamiast kupować nową na wyjazd.

\begin{zasada}[Policz przed planowaniem]
Nie planuj okna, zanim wiesz, skąd bierze się czas. Rachunek z dnia audytu jest tańszy niż
plan, który po tygodniu okazuje się niewykonalny, i daje konkretną liczbę zamiast poczucia winy.
\end{zasada}

\begin{uwaga}[Liczby z przykładów]
Dane w tabeli i w tekście są wyliczone z przyjętych założeń: 250 słów na minutę, książka
o długości \val{typical.words} słów, 50 minut dojazdu. Nie są pomiarami kogokolwiek. Zastąp je
własnymi, a rachunek zostanie ten sam.
\end{uwaga}

\section{Czytanie rodzinne i tablica minut}

Czytanie w rodzinie zwykle nie zwiększa dostępnego czasu. Zmienia tylko to, z kim go dzielisz.
Wspólne okno, na przykład 20 minut wieczorem, w którym każdy czyta własną książkę w tym samym
pokoju, daje tygodniowo 140 minut, czyli tyle samo, co wcześniej zabrało jedno przeglądanie
aplikacji. Nie trzeba udawać, że to dodatkowe godziny. Trzeba zdecydować, że są wspólne.

Czy dziecko widzi dorosłego czytającego obok, może mieć znaczenie dla jego nawyków, ale ta
korzyść nie jest w tym rozdziale przedmiotem rachunku. Każde twierdzenie o skutkach wspólnego czytania
w domu wymagałoby \verify{wspólne czytanie w domu wiąże się z lepszymi wynikami czytelniczymi dzieci}
sprawdzenia w literaturze, zanim trafiłoby do argumentacji. W tym
rozdziale mówimy tylko o czasie: czy dwie osoby mogą zająć ten sam wieczór książką.

Na koniec tablica minut, którą warto mieć przy lodówce. Dla dwunastu książek rocznie potrzeba
\val{need.12} minuty dziennie. Dla dwudziestu czterech \val{need.24} minut. Ten rozdział nie
mówi, że musisz czytać każdego dnia przez pół godziny. Mówi, że rachunek jest prosty i że
liczby z tablicy są mniejsze, niż wiele osób zakłada, zanim spojrzy na własny dzień.

\begin{przyklad}[Przykład dnia]
Przykładowa osoba, która chce przeczytać dwanaście książek w roku, zapisała przez tydzień swój dzień.
Okazało się, że dojazd daje jej 40 minut dziennie, a w kolejkach i przy oczekiwaniu na
spotkania przechodzi kolejne 15. Okno poranne zajmuje jej 20 minut, zanim zacznie pracę.
Mikroczas wystarcza na około 18 minut dziennie przy założeniu, że czyta w autobusie, a w
kolejkach tylko wtedy, gdy ma książkę pod ręką. Razem, czyli 18 minut mikroczasu i 20 minut okna, to 38 minut, ponad trzykrotność
\val{need.12} minut potrzebnych na dwanaście książek. Wieczorem przestała scrollować w łóżku,
a telefon ładuje w kuchni. Zmiana ta nie wymagała zakupu czegokolwiek poza budzikiem.
\end{przyklad}

\begin{falle}
\textbf{\glqq Czytam w przerwach, więc wystarczy\grqq{}.} Przerwa jest prawdziwym oknem tylko wtedy,
gdy ma stałe miejsce w dniu. Jeśli za każdym razem decydujesz, czy ją wykorzystać, decyzję podejmuje
zmęczenie. Przerwa, która zawsze kończy się telefonem, nie jest mikroczasem, tylko kolejnym
przeglądaniem.
\end{falle}

\begin{cwiczenia}
\exitem{Zapisz przez siedem dni trzy bloki swojego dnia dłuższe niż pięć minut, z liczbą minut
i tym, czy był w nich ekran.}{Trzy wiersze w formacie: blok, minuty, tak lub nie. Dobrze, gdy
jeden z bloków jest dojazdem, a jeden przerwą. Sprawdź, czy suma minut przekracza
\val{need.12}, bo wtedy mikroczas sam wystarcza na dwanaście książek.}
\exitem{Policz mikroczas z dnia: ile minut dziennie zostaje po odjęciu bloków, w których czytanie
jest niemożliwe, i ile z tego wykorzystasz realnie.}{Suma bloków możliwych do czytania, pomnożona
przez przyjętą część (w przykładzie jedna trzecia). Dla 88 minut wychodzi około 29 minut, co
wystarcza na dwadzieścia cztery książki, bo \val{need.24} minut. Wynik ma być liczbą, którą
możesz uzasadnić, a nie życzeniem.}
\exitem{Ile godzin rocznie daje codzienne okno 20 minut?}{$20 \times 365 = 7\,300$ minut, czyli
około 121{,}7 godziny.}
\exitem{Scrollowanie zajmuje 45 minut dziennie. Ile dób w roku oznacza to, licząc dobę jako 24 godziny?}{$273{,}75 / 24 \approx 11{,}4$, czyli około jedenastu i pół doby.}
\exitem{Podróż trwa sześć godzin, z czego czytasz 60\%. Ile minut zostaje na lekturę?}{$360 \times 0{,}6 = 216$ minut, czyli 3 godziny 36 minut.}
\exitem{Aplikacja zabiera 20 minut dziennie. Ile książek tygodniowo mogłeś przeczytać w tym czasie, przy tempie 250 słów na minutę i książce 80\,000 słów?}{$20 \times 7 \times 250 = 35\,000$ słów, czyli około $0{,}44$ książki tygodniowo. Liczba zależy od przyjętego tempa, więc sprawdź ją na własnym.}
\exitem{Dlaczego ładowarka poza sypialnią jest zmianą miejsca, a nie charakteru?}{Ponieważ nie wymaga decyzji w chwili pokusy. Przeglądanie zaczyna się samo wtedy, gdy telefon leży przy łóżku, a jego odsunięcie zmienia to, co jest wygodne, nie to, czego się chce.}
\end{cwiczenia}

\zadanie{Przez siedem dni zapisuj bloki dłuższe niż pięć minut i zsumuj je na koniec tygodnia.
Następnie zarezerwuj w kalendarzu jedno stałe okno, na przykład 20 minut rano, na kolejne siedem dni.
Miara sukcesu: tabela z co najmniej trzema blokami i wyliczoną sumą mikroczasu oraz okno
zrealizowane w co najmniej pięciu z siedmiu dni.}
